\section{Audit findings, corrections, and failed shortcuts}
\label{ado:sec:audit}
The source review was targeted, not a line-by-line audit of the entire
canonical manuscript. It covered the signed-kernel,
real-order, positive-double-kernel, subcritical, local-radius, and
research-status sections. No substantive counterexample to the reviewed
depth-two kernel and zero theorems was found. The following two
editorial corrections are concrete; the later items are scope safeguards
for this continuation rather than allegations that the manuscript makes
those mistakes.

\subsection{Two notation corrections in the transport proof}
In \path{chapters/05-real-positive-kernel.tex}, at the pinned snapshot,
the proof of the Hurwitz--polylogarithm transport formula says to
specialize the regularized Hurwitz identity ``with $q=n$'', and then
invokes boundedness of ``$h$''~\cite{repoPositive}. The displayed
Hurwitz formula uses $x$ as its real argument, and the regularized
kernel is named $q(t)$. The intended text is therefore
\[
 \text{``with $x=n$''},\qquad
 \text{``boundedness of $q$''}.
\]
The theorem and the coefficient calculation are unchanged. The file
\path{integration/notation_corrections.patch} contains precisely
these two replacements, at source lines 102 and 104. It is an optional
proposed patch; it has not been applied to the remote repository.

\subsection{Positive higher-order resolvents need not be zero-free}
The following exact example shows why the oscillation step in this
article cannot be replaced by a generic positivity slogan:
\begin{equation}\label{ado:eq:badpositive}
 R(z)=z^4\left(1+\frac1{(1-z)^4}\right)
     =z^4\int_{[0,1]}\frac{\dd(\delta_0+\delta_1)(u)}{(1-uz)^4}.
\end{equation}
The measure is positive, yet
\[
 z_0=1-e^{-\ii\pi/4}
     =1-\frac1{\sqrt2}+\frac{\ii}{\sqrt2}
\]
is a nonzero zero of $R$ and
$|z_0|^2=2-\sqrt2<1$. Indeed $(1-z_0)^4=-1$. Thus a positive
generalized Stieltjes kernel of order four does not imply even disk
nonvanishing. Our ordinary positive transform appears only
\emph{after} the sign-change polynomial has been identified and the
initial moments have been canceled.

\subsection{The unit-radius endpoint cannot be assigned an arbitrary phase}
At $\rho=1$, the point $z=1$ lies on the removed cut. A proof of
angular existence that simply sets $\Arg L(1)=0$ may silently assume
a finite endpoint value or a particular asymptotic. The proof in
Section~\ref{ado:sec:angular} instead uses the sector bound
$\Arg H<\beta_1$ to show that the unwrapped phase is below $\pi$
near the initial endpoint. This suffices even when $L$ diverges there.
No theorem about a limiting singular value is smuggled into the argument.

\subsection{The angular count does not extend automatically beyond radius one}
For the elementary depth-two example
$L_{1,1}(z)=\tfrac12\log^2(1-z)$, a zero of its imaginary part in
the upper half-plane requires
\[
 |1-z|=1.
\]
On $|z|=\rho$, this is $\cos\theta=\rho/2$. For $\rho>2$ there
is no such open-upper-arc zero at all. In our proof, the map
$t(u)=(1+\rho^2u^2)/(2\rho u)$ ceases to be monotone on $(0,1)$
when $\rho>1$, and the derivative used in the radial argument lemma
can also change sign. These are genuine losses of hypotheses.
Local continuation of a simple zero through $\rho=1$ must not be
confused with a theorem for all radii.

\subsection{Integer leading orders, normalized radii, and arithmetic depth}
The operator $\Aop^{s_j-1}$ was iterated an integer number of times.
Its elementary Rolle argument does not automatically apply to an
arbitrary fractional power. The manuscript itself proves that finite
signed moment representations can fail in part of the positive
real-order depth-two range~\cite{repoReal,repoSubcritical}. Therefore
one cannot state the present $L^1$ construction for all positive real
indices without a separate theorem or a different representation.

Likewise, strict decrease of $\theta_k(\rho)$ does not imply a chosen
sign for $\cos\theta_k(\rho)/\rho$. The original normalized-radius
conjecture and the fractional turning-point issues remain separate.
Finally, the degree $d-1$ in the compensation theorem is not a claim
about minimal multiple-polylogarithm depth at a numerical argument.
Formal relation modules, arithmetic period independence, and positive
transform compensation are different mathematical objects.

\subsection{The remaining arithmetic candidates retain their status}
The reviewed README and discovery chapter distinguish the proved $S_4$
identity, the unresolved $S_6$ and new $S_8$ candidates, and an older
$S_8$ vector that was rigorously rejected~\cite{repoStatus}. The
present analytic theorem does not furnish a relation certificate for
$S_6$ or $S_8$. Its Gaussian enclosures also do not establish arithmetic
independence or any conjectural equality merely by small residuals.
These status labels should remain unchanged when the package is
integrated.
